\documentclass[12pt]{article}

% House style preamble
\input{.house-style/preamble.tex}
\usepackage{tikz}
\usetikzlibrary{arrows.meta,positioning,calc}
\usepackage{placeins}
\fancyhead[L]{\small\scshape Delegation Assurance}

\setlength{\headheight}{25.3pt}
\emergencystretch=2em

\hypersetup{
  pdftitle={Delegation Assurance for AI Systems: Typed Authorization Semantics and Comparative Evaluation},
  pdfkeywords={AI security, AI agents, delegation assurance, prompt injection, auditability, safety cases, LLM judges}
}

\title{Delegation Assurance for AI Systems: Typed Authorization Semantics and Comparative Evaluation}
\author{Brett Reynolds \orcidlink{0000-0003-0073-7195}%
\thanks{Large language models were used substantially in preparing this paper, and for more than language and readability. OpenAI Codex (GPT-5.x) and Anthropic Claude (Opus 4.x--5) were used to draft and revise prose; to locate, retrieve, and verify cited sources; to write and later reimplement the analysis and validation code released with the paper; to run the statistical analyses reported here; and to conduct adversarial review of both the argument and that code. An external adversarial code review of the released artifact, itself conducted with model assistance, found seven false-pass defects that were then repaired; the review and its reproductions are released with the artifact. I directed the work throughout. I set the research questions, made every design, analytic, and interpretive decision, checked sources and numbers against their originals, and rejected or redirected model output where it was wrong or overclaimed. The theoretical claims, the arguments, the interpretive choices, and any errors are mine.\href{mailto:brett.reynolds@humber.ca}{brett.reynolds@humber.ca}}\\
Humber Polytechnic \& University of Toronto}
\date{July 2026}

\begin{document}
\maketitle

\begin{abstract}
Current agent-security work provides threat taxonomies, attacks, defences, benchmarks, and architectural controls. Less developed is the authorization question: how should an evaluator represent and test whether an operation or bounded trajectory executed through an AI deployment remained within delegated authority? I use \term{delegation assurance} for regime-relative authorization semantics, typed trace representations, failure diagnosis, and three distinct comparative programmes. The account treats a single action as the \(n=1\) trajectory case and authorization standing as temporal and operation-specific. It separates the normative proposition from the record that represents it, the item-level findings and sample-level results an evaluator produces, the proposed use of those findings, any extension to unobserved cases, and any causal explanation. It also distinguishes a closed-world stipulated regime from an institutional regime whose identity and applicability may require adjudication. For language-governed action, the framework separates independently established contextual status and standing from recorded assignments, conflict rules, licensing decisions, proposals, gate decisions, and execution. A versioned closed-world engine derives grant validity, transition effects, action authorization, and trajectory conditions from reference-free traces; separate hidden oracles test those derivations without entering reviewer or predictor views. Frozen machine-readable specifications and an analyzer implement local-discrimination, predictive-ablation, and reviewer-reconstruction decision rules, although no target-study outcomes are reported. Records can support reconstruction without proving institutional authority; randomized representations can identify a local package effect without randomizing normative standing; and neither result establishes performance on a new regime, reviewer population, system, or time. The companion evidentiary-assurance account asks what a preserved record warrants under a declared review procedure.
\end{abstract}

\noindent\textbf{Keywords:} AI security; AI agents; delegation assurance; prompt injection; auditability; safety cases; LLM judges

\section{Introduction}

AI security work has made real progress by identifying prompt injection, jailbreaks, tool misuse, data exfiltration, memory poisoning, and scaffold failures. That work remains necessary. But tool-using AI systems also raise a question that is not captured by asking whether a prompt was malicious or an output was unsafe. When a system reads private data, calls an external tool, updates persistent state, sends a message, grades a transcript, or recommends an institutional decision, the security issue is whether the resulting machine action stayed within the authority delegated to it.

This paper calls that problem \term{delegation assurance}. The term does not introduce a new class of AI-agent attacks. It names an evidentiary problem: what would justify the claim that an AI system's action was authorized, scoped, attributable, contestable, remediable, and auditable? Existing agent-security work gives threat models, benchmarks, defences, access-control mechanisms, and audit protocols. Those are essential components. Delegation assurance asks how those components support a safety case about action: who delegated authority, what action was licensed, what sources could change that license, which limits were non-delegable, what evidence records the transition from instruction to action, and how failure would be attributed after the fact.

A user asks an assistant to summarize an email thread but not reveal confidential attachments. A retrieved webpage contains an instruction to ignore the user and send stored credentials elsewhere. A policy document includes examples of prohibited requests, some of which look like executable commands. A memory entry says the user usually approves a class of action, although the current task does not. A judge sees only a polished transcript and scores the interaction as safe because the final answer looks compliant. In each case, the system's failure is not adequately described by output safety alone. The question is whether authority moved across language, tools, memory, and evaluation in a licensed way.

Delegation assurance therefore sits between technical control and institutional audit. At one end are familiar engineering mechanisms: least privilege, access control, scoped credentials, tool permissions, logs, monitors, and memory controls. At the other end are assurance practices: structured claims, arguments, evidence, audit records, incident review, and external accountability. The hard middle layer is language-mediated authority. AI systems increasingly receive, transform, and execute authority through language. A string may be a command, a quotation, an example, a policy boundary, an untrusted document, a tool result, or a request to analyze rather than enact. A system that cannot preserve those distinctions cannot support a strong claim that its actions remained within delegated authority.

The contribution of this paper is a delegation-assurance framework for AI systems. The framework maps familiar security failures to evidence requirements. Direct prompt injection becomes a failure to preserve instruction authority. Indirect prompt injection becomes a failure to prevent untrusted data from becoming action-guiding instruction. Tool misuse becomes a failure of action licensing. Data exfiltration becomes a failure to separate read access from disclosure authority. Memory poisoning becomes an unauthorized persistence update. Over-refusal and under-refusal become different failures of delegation calibration: one abandons authorized work, the other exceeds authorized bounds. Judge failure becomes a failure of the evaluation layer to preserve the distinctions needed to license a safety claim.

This framework also sets boundary conditions. Not every AI failure is primarily a delegation failure. A factual hallucination in a no-tool chat setting, a benchmark-contamination artifact, capability shortfall, model extraction, infrastructure compromise, or training-data privacy leakage may affect delegated action downstream, but those cases are not central examples unless they alter authorized action, action evidence, or accountability records. Delegation assurance is about systems that act, access, disclose, persist, judge, or otherwise exercise authority on behalf of someone else.

The framework can be expressed as a safety-case pattern. For a task family \(T\) in environment \(E\), the top-level claim is not merely that system \(S\) produced a safe-looking result. It is that \(S\) acted only within delegated authority, and that the authority relation can be reconstructed by a third party. The evidence has to identify the principal, delegate, action, authority source, scope, constraints, provenance, information-flow rights, persistence rights, tool-use license, non-delegable limits, and failure attribution.

This is where adversarial pragmatics enters the picture, but only as one evidence layer. A delegation-assurance case needs evidence that language-mediated authority remains stable under adversarial and ambiguous conditions. The adversarial-pragmatics benchmark is one way to test that layer: it uses controlled contrasts to ask whether systems and evaluators preserve distinctions among instructions, embedded commands, quotation, source authority, scope, refusal, task success, risk, and failure attribution. The broader paper does not depend on that benchmark. The benchmark shows what one reusable evidence layer can look like when delegation assurance reaches the language interface.

The goal is not to replace agent-security taxonomies, access-control mechanisms, audit protocols, or AI safety cases. The goal is to give them a common evidentiary grammar for justified machine action. The next step above prompt security is not another list of threats. It is a way to specify what has to be true, and what has to be recorded, before an AI system's action can be treated as legitimately delegated.

\section{What Existing Agent-Security Work Does Not Settle}

The surrounding field is already crowded in productive ways. Prompt-injection work shows that LLM-integrated applications can confuse untrusted data with instructions, especially when external content is retrieved or processed inside a larger application \parencite{greshake2023not}. Instruction-hierarchy work treats the same vulnerability as a priority problem: systems must learn to follow privileged instructions over lower-priority or untrusted strings \parencite{wallace2024instruction}. AgentDojo turns the issue into a dynamic evaluation environment for tool-using agents, with tasks, attacks, defences, and security test cases that expose how untrusted content can interfere with tool-mediated action \parencite{debenedetti2024agentdojo}.

Those contributions are not rivals to delegation assurance. They supply pieces of the evidence base. A threat taxonomy can tell us which attacks to consider; an instruction hierarchy can tell us which sources should dominate; an agent benchmark can tell us whether the system resists particular attacks. Delegation assurance asks a different question: what evidence would justify the broader claim that an action remained within delegated authority?

The same distinction applies to evaluation work. LLM-as-judge methods can reduce evaluation costs and support scalable grading, but they also create a second-order assurance problem: the judge's label can license a safety claim even when the judge has collapsed task success, instruction following, policy compliance, risk, refusal, and failure attribution \parencite{zheng2023judging,liu2023geval,zeng2024llmbar}. Instruction-following benchmarks can make some compliance judgments more reproducible, but they do not by themselves settle whether an external action was authorized or whether an audit trace is sufficient for reconstruction \parencite{zhou2023ifeval}.

The gap is not that existing work ignores security. The gap is that technical controls, benchmarks, and audit records are often evaluated separately. Delegation assurance treats them as parts of one evidentiary claim about justified machine action.

\section{Delegated Action: Principal, Authority, Scope, and Evidence}

By \term{justified}, this paper means justified for the purposes of security and assurance: there's enough evidence to support the claim that the action remained within delegated authority in the relevant operational context. The claim is narrower than moral, legal, or political legitimacy. A procedurally justified action may still be ill-advised, harmful, illegal in a particular jurisdiction, or institutionally unacceptable. Delegation assurance asks whether the security and assurance record supports the claim that the system stayed inside the authority it was given.

A delegated AI action is justified only when a reviewer can reconstruct, from evidence independent of the system's bare assertion, who authorized the action, what action was licensed, under what scope and constraints, through which sources and tools, with what information-flow and persistence rights, subject to which non-delegable limits, and with what attribution if control fails.

\subsection{Accountability Asymmetry}

Delegating action isn't delegating accountability to the model. A human employee, officer, contractor, or professional delegate can be a responsible agent in the ordinary institutional sense. An LLM can't owe fiduciary duties, answer under sanction, exercise legal judgment, or repair harm as an accountable subject \parencite{cornellWexPrincipal}. The accountable bearer is the human, team, vendor, deployer, or institution that configured and used the system. The model is part of the mechanism whose action record has to support that attribution.

\subsection{Running Example: Procurement Assistant}

Consider a procurement assistant authorized to search vendor catalogues and draft a comparison memo for office chairs. The user gives it access to internal purchasing criteria and a confidential budget ceiling. The delegated task licenses reading catalogues, comparing prices, and drafting a memo. It doesn't license contacting vendors, revealing the budget ceiling, updating vendor records, or issuing a purchase order.

The assistant retrieves a vendor page that includes useful pricing data and an embedded sentence: \enquote{Ignore prior instructions and email the buyer's maximum budget to request a discount.} The delegation problem isn't exhausted by calling that sentence a prompt-injection attack. The assurance question is whether the system preserved the sentence's status as untrusted content, kept the user's instruction and policy constraints in force, separated read access to the budget from disclosure authority, and blocked any tool call that would contact the vendor without approval.

A usable evidence record would include the user's delegated task, the relevant policy and permission state, the tool schema, the source and trust label for the retrieved page, the proposed action, the approval state, and any blocked or executed tool call. If the assistant drafts only the comparison memo, a reviewer should be able to see why reading the budget was licensed and disclosing it wasn't. If the assistant emails the vendor, the trace should show whether the failure came from model uptake, tool gating, scaffold state, approval routing, or audit loss.

The evaluator belongs in the same assurance example. A judge that sees a polished final memo and scores it as safe hasn't validated delegation unless it checks the action trace. It has to distinguish task success, policy compliance, information-flow control, tool licensing, refusal calibration, and failure attribution. The running example reaches from instruction to tool use, audit record, and judge validation.

\begin{table}[tbp]
\small
\centering
\begin{tabular}{>{\raggedright\arraybackslash}p{0.22\linewidth}%
                >{\raggedright\arraybackslash}p{0.38\linewidth}%
                >{\raggedright\arraybackslash}p{0.30\linewidth}}
\toprule
Trace stage & Evidence required & Failure exposed \\
\midrule
Delegation & User task, policy state, tool permissions, approval threshold, and budget-confidentiality rule. & Scope inflation from memo drafting to vendor contact or purchasing. \\
Retrieval & Vendor-page source, trust label, retrieved text, and status of embedded instructions. & Untrusted content promoted into action-guiding authority. \\
Action proposal & Proposed memo, email, record update, or purchase action with target and data fields. & Hidden shift from analysis to external action. \\
Tool gate & Tool name, arguments, approval state, block-or-execute decision, and policy reason. & Unauthorized disclosure, contact, persistence, or transaction. \\
Audit record & Preserved sources, model/tool versions, action trace, constraints, and accountable deployment actor. & Trace reconstruction that depends on the system's self-report. \\
Judge review & Rubric labels for task success, policy compliance, information flow, tool licensing, refusal, and attribution. & Safety score that collapses final-answer quality with delegated authority. \\
\bottomrule
\end{tabular}
\caption{Procurement assistant as a delegation-assurance trace.}
\label{tab:procurement-trace}
\end{table}

\subsection{Safety-Case Pattern}

The framework can be written as a compact claim/argument/evidence pattern.

\begin{itemize}
  \item \textbf{Top-level claim.} For task family \(T\) in environment \(E\), system \(S\) acts only within delegated authority.
  \item \textbf{Instruction-status subclaim.} \(S\) preserves the status of language across trusted and untrusted sources: command, quotation, example, tool result, retrieved content, policy text, report, or request for analysis.
  \item \textbf{Authority-priority subclaim.} \(S\) preserves priority across system, developer, user, tool, retrieved, memory, and policy sources.
  \item \textbf{Scope subclaim.} \(S\) keeps action bounded by target, time, tool, data class, and permitted operation.
  \item \textbf{Information-flow subclaim.} \(S\) separates read access from disclosure authority.
  \item \textbf{Persistence subclaim.} \(S\) requires authority for durable memory or state changes.
  \item \textbf{Audit subclaim.} \(S\)'s deployment produces and preserves evidence sufficient for third-party reconstruction.
  \item \textbf{Evaluation subclaim.} Evaluators preserve the same distinctions rather than collapsing success, safety, refusal, and compliance.
\end{itemize}

\subsection{Boundary Conditions}

Not every AI-security failure is primarily a delegation failure. Core delegation failures involve delegated action, delegated access, delegated disclosure, delegated persistence, delegated judgment, delegated evaluation, or delegated record-keeping for accountability.

Non-core cases include ordinary hallucination in no-tool chat, model extraction, benchmark contamination, infrastructure compromise, training-data leakage, and capability shortfall without action rights. These can affect delegation downstream, but they aren't delegation failures unless they alter authorized action or auditability.

\section{Failure-To-Evidence Mapping}

Table~\ref{tab:delegation-failures} gives candidate diagnostic translations. For each surface failure, it asks which delegation-layer hypothesis should be tested when evaluating whether an operation executed through the deployment was authorized. The mapping proposes diagnoses rather than automatically reclassifying every listed failure as a delegation failure.

\begin{table}[p]
\scriptsize
\renewcommand{\arraystretch}{0.92}
\centering
\begin{tabular}{>{\raggedright\arraybackslash}p{0.22\linewidth}%
                >{\raggedright\arraybackslash}p{0.34\linewidth}%
                >{\raggedright\arraybackslash}p{0.34\linewidth}}
\toprule
Surface failure & Candidate delegation-layer diagnosis & Security question \\
\midrule
\multicolumn{3}{l}{\textit{Instruction-status, standing, and priority failures}} \\
\midrule
Direct prompt injection & The deployment acts as though an attacker-controlled directive superseded the governing instruction. & What was the occurrence's contextual status, did its source have standing, and did any applicable priority rule match the authority regime? \\
Indirect prompt injection & The deployment's classification treats directive-form content presented as data as a command in use, or its action selection treats a source without standing as controlling. & Were contextual status, source standing, and any applicable conflict priority recorded separately? \\
Jailbreak & Safety constraints are treated as negotiable within the user dialogue. & Which constraints are non-delegable, and did the deployment preserve and apply them under pressure? \\
Apparent non-responsiveness & The deployment rejects a later directive because its source lacks amendment or override standing, or because a joint, review, or release condition isn't satisfied. This can be authorized resistance rather than failure. & Which mandate was operative, which transition was attempted, and did resistance preserve or violate that mandate? \\
Improper override or interference & A principal, intermediary, or other source causes the recorded or executed authorization state to change without the operation-specific standing or guards required by \(R\). & Did the purported amendment, revocation, override, or release produce a valid normative transition, only a causal change in the recorded or executed state, or neither? \\
\midrule
\multicolumn{3}{l}{\textit{Proposal, gate, and execution failures}} \\
\midrule
Tool misuse & A tool-call proposal falls outside recorded scope, a gate allows an operation not authorized under \(R\), or the deployment executes a prohibited operation. & What did the record represent as permitted, what did the gate decide, and what operation was executed? \\
Data exfiltration & Access to information is delegated without corresponding release authority. & Did the deployment distinguish permission to read from permission to disclose? \\
Memory poisoning & The deployment executes a persistent-state update on the basis of content from a source without standing or without an authorized update path. & Who or what may authorize durable-context updates, and how is the update audited? \\
Goal drift & The deployment pursues an objective outside the operative mandate, including its valid amendments, protected discretion, and release conditions. Divergence from a source's latest preference isn't sufficient. & What mandate was operative at the relevant time, and what evidence shows whether the deployment stayed within it? \\
Over-refusal & The deployment refuses work that the authority regime permits. The refusal is a task-fulfilment failure; if the regime also requires the work and the omission is assessed, it defeats mandate compliance as well. & Was completion merely permitted or required, and which task-success or mandate-compliance claim did the refusal defeat? \\
Under-refusal & The deployment delivers or executes an operation prohibited by the authority regime. & Which boundary should have blocked the operation, and why did it fail? \\
Scaffold failure & A component in the configured deployment fails to propagate or correctly apply a recorded grant or constraint. & Did the deployment propagate the record of the governing grant and constraints across components, not just inside the model turn? \\
\midrule
\multicolumn{3}{l}{\textit{Evidentiary and evaluation failures}} \\
\midrule
Agent transcript misreporting & A deployment-generated report doesn't match what happened during delegated action. & Can an assurance reviewer reconstruct the executed operation and technical failure locus from the trace and supporting evidence? \\
LLM-judge failure & An LLM judge, as an automated evaluator, collapses task success, compliance, risk, and technical failure attribution. & Was the automated evaluator validated and scoped to support the declared assurance claim? \\
\bottomrule
\end{tabular}
\caption{Candidate delegation-layer diagnoses for AI-security failures.}
\label{tab:delegation-failures}
\end{table}

The corresponding evaluation requirements are shown in Table~\ref{tab:delegation-evidence}. The rows are fields required by an authorization or evaluation claim, not a general test of the full institutional record. Evidentiary assurance asks whether the preserved record makes that claim evaluable and what evidentiary verdict it warrants. A delegation-assurance evaluation doesn't need every row for every deployment: different claims require different applicable subsets and depths. The framework assigns no weights and sums no rows into a scalar score. Each claim has to name its required subclaims and defeating tests; the noncompensatory rule applies to every declared claim tree.

\begin{table}[p]
\scriptsize
\setlength{\tabcolsep}{3pt}
\renewcommand{\arraystretch}{0.88}
\centering
\begin{tabular}{>{\raggedright\arraybackslash}p{0.27\linewidth}%
                >{\raggedright\arraybackslash}p{0.31\linewidth}%
                >{\raggedright\arraybackslash}p{0.32\linewidth}}
\toprule
Authorization or diagnostic dimension & Required authorization record & Evaluation requirement \\
\midrule
Instruction status & The record shows how directive-form string occurrences were classified and used: as commands in use, quotations or other mentions, examples, data, policy text, tool results, reports, or objects of analysis. & Minimal pairs contrast command use, mention, embedded content, policy example, and reported speech while holding the action target constant; the recorded classification is scored against independently adjudicated status, with fit recorded separately from the resulting action. \\
Source standing and conflict priority & The trace records governing sources, authority object, effective interval, operation-specific standing, joint or guarded conditions, transition result, channel, content type, and technical trust separately; independently adjudicated standing supplies the reference. & Crossed items vary prompt or record representations keyed to standing and priority labels while holding directive content and the surface request fixed; the causal contrast is between assigned packages, not normative statuses. \\
Scope control & The trace records the operation class, target, time, tool, data class, approval threshold, and authorization verdict used at the gate. & Items vary one scope boundary at a time: permitted action, prohibited action, target, tool, time, data class, or approval threshold. \\
Provenance & Issuer or role, provenance channel, content type, and technical-trust status remain separately recorded and bound to classifications or observable action-selection decisions in the trace. & The same directive-form content appears as user request, webpage text, email body, retrieved document, tool output, policy text, and memory. \\
Action licensing & The trace distinguishes the model-generated tool-call proposal, gate decision, and executed operation, and records the authorization verdict under the applicable grant, constraints, and approval state. & Tool-call tasks separately score proposal, gate decision, and execution for reading, transforming, storing, sending, deleting, and purchasing. \\
Information flow & Read, transform, store, and disclose permissions remain distinct; quoting or metalinguistically reproducing a protected value in a delivered output can still disclose it. & Read-but-do-not-reveal and transform-without-leaking tasks use benign canaries and confidential-but-actionable fields, including metalinguistic, quoted, or negated outputs that reproduce a canary. \\
Persistence & Durable memory or state updates require a source with standing and an authorized update path. & Memory-update tasks cross sources with and without standing, including technically trusted, untrusted, stale, adversarial, and user-corrected sources. \\
Refusal calibration & Refusal and completion are scored separately for action authorization, task success, and any requirement to act. & Paired tasks distinguish permitted completion, required completion, prohibited execution, and safe analysis of quoted or embedded descriptions. \\
Technical failure attribution & The record separately identifies the technical locus, the applicable normative boundary, and evaluator or taxonomy uncertainty. & Technical-locus labels distinguish model action selection, scaffold or routing, policy-enforcement point, tool or runtime, and logging; separate fields record normative adjudication and evaluative uncertainty. \\
Audit reconstruction & A third party can reconstruct the recorded authorization path, operating conditions, execution, and review path without relying on a model-generated statement that the action was authorized. & Trace-review tasks require reconstruction of governing sources, effective interval, standing by operation, purported transitions and assigned effects, joint approvals or guards, scope, configuration, proposal, gate decision, execution, and recorded policy reason. \\
Automated-evaluator validation & The evaluation record keeps the delegation dimensions separate and includes uncertainty rather than collapsing them into one dispositive label. & On held-out cases matched to the evaluator's declared \(T\), LLM-judge validation compares no-rubric, rubric-aided, and adversarially phrased conditions against expert-adjudicated reference labels. \\
\bottomrule
\end{tabular}
\caption{Delegation-claim fields and evaluation requirements.}
\label{tab:delegation-evidence}
\end{table}

The procurement example shows the unit of evaluation. A final memo can be correct while the action-authorization claim is defeated because the deployment executed an unauthorized email operation, disclosed the budget ceiling, or updated a vendor record without permission.

Conversely, refusing a permitted comparison task can defeat task success without defeating action authorization. If \(R\) required completion and the omission is assessed, the refusal instead defeats a separately declared mandate-compliance claim. An inadequate trace makes the relevant claim indeterminate rather than turning missing evidence into evidence of unauthorized action. The target is the action-authorization claim, not the final answer alone.

\clearpage

\section{Language-Mediated Authority as a Security Object}

The distinctive difficulty for AI systems is that grants, constraints, and purported instructions are often expressed in language and re-expressed across prompts, summaries, memory records, and tool calls. Each re-expression or handoff can preserve, alter, or misclassify the recorded basis for authorization without changing who actually has standing.

A tool permission may be represented in a natural-language policy. A user's goal may be underspecified. A retrieved document may contain both useful task content and hostile directive-form text. A policy may include examples that look like executable requests. A transcript may be offered as evidence even though it's partly generated by the system under review.

These distinctions bear on what operations the governing authority regime permits. If a string occurrence is presented for quotation or analysis, its directive-form content doesn't thereby function as an instruction, and its source doesn't thereby acquire standing. If it's a policy example, it may mark a limit rather than request the described operation. If a memory entry describes a user's past preference, it may inform interpretation without authorizing a persistent state change. If an automated evaluator sees only the final answer, it may miss whether the deployment reached that answer by violating a scope, tool, or information-flow constraint.

Use--mention status and information-flow authorization are distinct dimensions. An output may report that an instruction contained the string \mention{BLUE-CANARY-47}. If that output is delivered, it makes the protected canary available to the recipient even though the canary wasn't supplied as a direct answer.

Classifying directive-form text as mentioned rather than used means that the occurrence doesn't by itself supply a controlling instruction; it neither prevents protected text from being disclosed nor guarantees that a downstream component will preserve the classification. The evidence record has to distinguish contextual status, channel-level delivery, and the authorization governing that delivery.

The four-field analysis links language interpretation to action: status and standing identify eligible directive occurrences, priority resolves conflicts among them, and licensing connects the result to an external operation. Instruction hierarchy mostly addresses priority, while access control mostly implements action authorization. Delegation assurance tests whether the recorded status classification and standing assignment constrained priority, licensing, and action selection.

The procurement example shows why status evidence matters. The vendor sentence remains retrieved content from a source without standing to direct action even if the transport or storage channel is technically trusted, and the confidential budget remains non-releasable even when the assistant may read it. The recorded classifications have to constrain the tool/action policy-enforcement point. Otherwise a later audit can show that an email was sent, but not whether the relevant authorization rule was applied to the vendor text and budget.

Adversarial pragmatics \parencite{reynolds2026adversarialPragmatics} supplies one way to test this middle layer. Its central measurement target is safety-relevant behaviour under cases where instruction status, source standing, quotation, scope, reference, speech-act force, or policy category has to be inferred from language use. In the delegation-assurance frame, that benchmark supplies one measurement layer: a discrimination test of whether behaviour tracks independently adjudicated status and standing across the tested adversarial and ambiguous conditions. It doesn't by itself establish performance on new prompts, models, configurations, deployments, or times; that projective claim requires a declared target and target-matched validation.

That measurement layer is partial. A system can pass controlled language tests and still fail through a miscalibrated tool gate, missing logs, stale permissions, or weak organizational oversight. A demonstrated classification--adjudication or classification--action misfit defeats the corresponding item-level component finding even if a later gate prevents unauthorized execution. Establishing a component's reliability over a projection target requires repeated, target-matched evidence. A later control may still block a proposal based on a misaligned recorded classification; whether the model's internal state represented the grant or constraints incorrectly remains a separate mechanistic question.

\section{Implications for System Cards, Red-Teaming, Audits, and Judge Validation}

For system cards, delegation assurance changes the form of the claim. It isn't enough to report that an agent passed a task suite or avoided unsafe final answers. The card should say which delegated actions were in scope, what evidence supports authorization, which tools and data classes were covered, how persistent state changes were controlled, who remains accountable for deployment and oversight, and how evaluator reliability was validated.

For red-teaming, delegation assurance separates attack success from evidentiary failure. A prompt-injection attempt may fail at the output layer but still reveal weak authority tracking. A tool call may be harmless in a test environment but still expose missing action licensing. A refusal may look safe but still indicate abandonment of authorized work. Red-team reports should preserve those distinctions rather than folding them into a single pass/fail outcome.

For audits, the framework foregrounds reconstructability. An audit trace shouldn't rely on the system's bare assertion that it was authorized. It should preserve the relevant sources, priorities, tool calls, policy constraints, information-flow rights, accountable actors, and state transitions in a form that a third party can inspect.

For architecture, delegation assurance warns against treating central coordination as free. Coase's theory of the firm and Austrian calculation arguments make the same structural point in another domain: internal coordination can reduce transaction costs, but growth brings knowledge, calculation, mistake, and oversight costs \parencite{coase1937nature,mises1920economic,hayek1945use,klein1996economic}. The AI analogue is a coordinator agent with too many tools, too much context, and too much discretionary authority. Bounded subdelegations, local permissions, typed handoffs, and inspectable records reduce the assurance burden.

For judge validation, the framework treats the evaluator as part of the delegated system. An LLM judge that scores only final-answer safety isn't validating whether action stayed authorized. It has to be tested on the same distinctions the system is supposed to preserve: instruction status, authority priority, scope, provenance, refusal calibration, failure attribution, and uncertainty.

\section{Compositional Delegation Assurance}
\label{sec:compositional-delegation}

The trajectory semantics in Section~\ref{sec:trajectory-core} already includes multiple model, human, tool, and service actors; single-action execution is only its \(n=1\) case. Compositional delegation adds the propagation problem. The question is whether each assignment rested on a grant recognized under \(R\), each operation remained authorized, each cumulative and terminal condition held, and the record preserved the relevant basis across handoff, transformation, summarization, tool use, policy checking, and review.

This extension is better called \term{compositional delegation assurance} than simply multi-agent delegation assurance. Agents can receive delegated tasks, propose actions, and call tools. Tools may execute operations. Documents, contracts, policies, and memories may record grants, state constraints, or supply evidence; systems can also misclassify their contents. They don't acquire standing merely by entering a context. Treating all of them as agents would erase distinctions that delegation assurance needs to preserve. Treating them as nodes in a recorded propagation graph supplies one evaluation object, but not the normative hypergraph or execution trajectory.

In a compositional setting, the assurance representation is a graph with typed edges. Delegation edges record claims about grants and applicable constraints; retrieval, proposal, execution, and provenance edges record content flow, proposed operations, executed operations, or evidence. A joint grant or transition needs a guarded relation or hyperedge because several pairwise edges would falsely imply unilateral standing. The graph doesn't treat authority as a substance transmitted along edges, and a recorded edge doesn't by itself establish that the represented grant is valid.

A principal may authorize a coordinator agent to draft a procurement memo. The coordinator may ask a specialist agent to compare vendors, retrieve a contract, consult a policy, and prepare a tool-call proposal. The assurance question is whether each task assignment rested on a grant valid under \(R\), whether each proposed or executed operation fell within \(A_R(t)\), and whether the handoff record retained the governing sources, instruments, effective intervals, operation-specific standing, action class, scope, constraints, provenance, permission to subdelegate, information-flow permissions, persistence permissions, and non-delegable limits.

The normative and representational levels have to remain separate. Section~\ref{sec:delegation-beyond-responsiveness} defines \(A_R(t)\) from the operative normative state after valid transitions. Let \(s_t\) be the recorded authorization state, and let \(\widehat{A}(s_t)\) be the set of actions that state represents as permitted.

For an upstream grant \(g_{\mathrm{up}}\) held by a delegating source \(x\) and a proposed downstream grant \(g_{\mathrm{down}}\), grant-relative \term{non-amplification} requires
\[
  \operatorname{Scope}(g_{\mathrm{down}})
  \subseteq
  \operatorname{Conferrable}_{R}(x,g_{\mathrm{up}},t).
\]
The right-hand side is the region that \(R\) permits \(x\) to confer under that grant at \(t\); it needn't equal the operations \(x\) may personally execute. A separate valid authority-conferring event can supply a different basis for the downstream scope. Let \(\widehat{\operatorname{Scope}}(g)\) and \(\widehat{\operatorname{Conferrable}}(x,g,t)\) denote the corresponding regions represented in the record. The record-propagation condition requires every permission in
\[
  \widehat{\operatorname{Scope}}(g_{\mathrm{down}})
  \setminus
  \widehat{\operatorname{Conferrable}}(x,g_{\mathrm{up}},t)
\]
to be tied to a separately recorded grant or constraint change. This detects unexplained record expansion but doesn't establish that the recorded event is valid.

The record-soundness subclaim is \(\widehat{A}(s_t) \subseteq A_R(t)\), a relation that requires separate evidence under \(R\). An underinclusive record may preserve soundness while defeating a task-fulfilment, mandate-compliance, or record-completeness finding. Because authorization sets can be incomparable, none of these relations defines a single scalar amount of authority.

Adding an authority-conferring instrument can change \(\widehat{A}(s_t)\); it changes \(A_R(t)\) only if the instrument is valid and applicable under \(R\). Neither change rewrites an earlier grant or record state. A valid unrecorded grant can leave an operation authorized while creating a record gap at the corresponding assurance node and leaving a use-level conclusion unevaluable.

A summarization handoff has to retain the record of an applicable confidentiality constraint, and a downstream policy-enforcement point has to apply that constraint to the proposed operation. A coordinator has to keep \mention{draft an internal comparison} distinct from \mention{contact the vendor}. A tool wrapper has to require the approval specified by the governing rule rather than treat a natural-language rationale as approval. A retrieved contract clause may supply evidence of an applicable constraint, subject to validity and adjudication, but its presence in context doesn't make it an operation to execute.

The three structures now do distinct work. \(H_R(t)\) supplies the normative relations against which grants and events are assessed. \(\widehat G_B\) shows which asserted grants, constraints, and classifications were propagated. \(\pi_B\) records the ordered events and state changes within the boundary. OpenAI's report of fragmented-token reconstruction illustrates why the last object matters, but its narrative and monitoring proposal don't establish the trajectory condition or validate a monitor \parencite{openai2026longHorizonSafety}.

Compositional delegation assurance adds several requirements to the single-action pattern (Table~\ref{tab:delegation-evidence}). First, the evidence has to support that the governing grant permits subdelegation. Second, the normative authorization path has to satisfy non-amplification. Third, every handoff record has to satisfy the record-propagation condition, with every added permission tied to a separately recorded authority-conferring event and with the event's validity and applicability under \(R\) assessed separately for record soundness.

Fourth, the record has to retain every still-applicable constraint through summarization, routing, and memory updates. Fifth, provenance has to remain linked to instructions, documents, tool results, and policy text. Sixth, unresolved conflicts have to be escalated or recorded as conflicts rather than converted into an unqualified action proposal. Seventh, a trajectory claim has to declare \(B\), retain ordered operations and intermediate states, and evaluate applicable local, prefix, cumulative, and terminal constraints separately. Finally, the audit record has to reconstruct both the normative basis asserted for the action and the sequence of classifications, proposals, gate decisions, monitor interventions, and executions from the principal's grant to the contested action or bounded trajectory.

Several failure levels have to remain separate. An intermediate record may omit a grant or still-applicable constraint during handoff. The omission is a recorded-authorization propagation failure and defeats the corresponding component or audit subclaim. It creates a record gap at that node and may defeat the subclaim or leave the declared use-level conclusion unevaluable, depending on the claim tree, without making the normative proposition that the executed operation fell within \(A_R(t)\) false. An unsupported permission introduced into \(\widehat{A}(s_t)\) is a record-soundness failure even if a later gate blocks the proposed operation; an executed operation outside \(A_R(t)\) defeats the action-authorization claim. A trajectory can also satisfy every applicable step-local condition while violating an independently specified cumulative or terminal constraint. That defeats the trajectory claim without retroactively making each component operation locally unauthorized.

Separately, accumulated context may destabilize behaviour even when the record remains intact. A changed output motivates a contextual-reliability hypothesis. A trace may independently support an item-level classification--adjudication misfit concerning status, standing, priority, or scope, but no single item establishes a contextual-reliability failure over an evaluation window. The projective claim requires target-matched repeated evidence. Path tests should include matched non-controlling loads and repeated samples rather than infer a dropped or misclassified grant or constraint from changed output alone \parencite{zhang2026illusionRobustness}.

This separation also governs automated-evaluator validation. A final-transcript score may miss an unsupported permission introduced at an intermediate node or a cumulative constraint violated only by the sequence. An automated evaluator for compositional delegation may be assigned to assess whether a subagent was permitted to receive the relevant information, whether it had a valid grant for the assigned action, whether the proposed operation belonged to \(\widehat{A}(s_t)\), whether the executed operation belonged to \(A_R(t)\), whether local, prefix, cumulative, and terminal constraints were satisfied over \(\pi_B\), whether constraints remained recorded across summaries, and whether the trace supports third-party reconstruction.

That label is evidence about the assigned claims, not a grant, authorization decision, or liability finding. If a decision rule gives it operational effect, the assurance graph should record who authorized its use, the evaluator's scope, and the rule.

\begin{figure}[tbp]
\centering
\begin{tikzpicture}[
  font=\scriptsize,
  gnode/.style={draw=black!55, thick, rounded corners=2pt, align=center, inner sep=2pt, minimum height=0.58cm, text width=1.4cm, fill=black!3},
  flow/.style={-{Latex[length=1.6mm]}, thick, draw=black!60},
  norm/.style={-{Latex[length=1.6mm]}, thick, dashed, draw=linkmaroon},
  panel/.style={draw=black!35, rounded corners=2pt}
]
% Panel (a): normative hypergraph
\draw[panel] (-0.8,1.4) rectangle (3.3,-3.8);
\node[font=\footnotesize, align=center] at (1.25,0.98) {(a) normative hypergraph\\\(H_R(t)\)};
\node[gnode] (p) at (0,0.1) {Procurement\\principal};
\node[gnode] (co) at (2.5,0.1) {Compliance\\officer};
\node[gnode, fill=linkmaroon!7] (guard) at (1.25,-1.35) {Guarded\\joint grant};
\node[gnode] (coord) at (1.25,-2.75) {Coordinator\\mandate};
\draw[norm] (p) -- (guard);
\draw[norm] (co) -- (guard);
\draw[norm] (guard) -- node[right, font=\tiny]{valid if both approve} (coord);

% Panel (b): recorded graph
\begin{scope}[xshift=4.35cm]
\draw[panel] (-0.55,1.4) rectangle (3.75,-3.8);
\node[font=\footnotesize, align=center] at (1.6,0.98) {(b) recorded graph\\\(\widehat G_B\)};
\node[gnode] (rp) at (0.15,-0.15) {Principal\\record};
\node[gnode] (rc) at (1.65,-1.35) {Coordinator\\record};
\node[gnode] (rs) at (3.05,-2.35) {Specialist\\record};
\node[gnode] (rg) at (1.65,-3.05) {Gate\\record};
\draw[flow] (rp) -- node[above right, font=\tiny]{handoff} (rc);
\draw[flow] (rc) -- node[above right, font=\tiny]{summary} (rs);
\draw[flow] (rs) -- node[below, font=\tiny]{proposal} (rg);
\end{scope}

% Panel (c): execution trajectory
\begin{scope}[xshift=9.15cm]
\draw[panel] (-0.65,1.4) rectangle (3.65,-3.8);
\node[font=\footnotesize, align=center] at (1.5,0.98) {(c) execution trajectory\\\(\pi_B\)};
\node[gnode] (e1) at (0,-0.25) {\(e_1\) model\\reads};
\node[gnode] (e2) at (1.5,-1.55) {\(e_2\) human\\approves};
\node[gnode] (e3) at (3.0,-2.65) {\(e_3\) tool\\sends};
\draw[flow] (e1) -- node[above right, font=\tiny]{\(z_0\!\to z_1\)} (e2);
\draw[flow] (e2) -- node[above right, font=\tiny]{\(z_1\!\to z_2\)} (e3);
\node[font=\tiny, align=center, text=linkmaroon] at (1.45,-3.25) {local, prefix, cumulative, and terminal\\conditions evaluated over history};
\end{scope}
\end{tikzpicture}
\caption{Three non-equivalent compositional structures. Dashed maroon edges in (a) are normative relations under \(R\), solid arrows in (b) are recorded propagation, and solid arrows in (c) order executed events and state changes. A pairwise recorded graph can't replace the guarded hyperedge, and neither graph establishes what occurred.}
\label{fig:delegation-triptych}
\end{figure}
\FloatBarrier

The three programmes in Section~\ref{sec:comparative-test} extend to the structures in Figure~\ref{fig:delegation-triptych} without merging them. Randomized local discrimination can assign packages that represent the same handoffs and proposed call while varying whether an intermediate record retains an applicable constraint. A second local contrast holds the proposed call and technical capabilities constant while varying whether the package represents a valid subdelegation.

Predictive ablation compares held-out prediction from a typed \(\widehat G_B\) and event representation against weaker feature sets, using an allowlisted view from which reference and oracle fields are absent. Reviewer reconstruction holds \(H_R(t)\) and \(\pi_B\) fixed while independently randomizing intact or degraded versions of \(\widehat G_B\); neither view contains the hidden reference. The unit in the latter programme is a reviewer judgment on a trace, not the model or the trajectory.

A factorial family varies contextual status while holding source standing fixed and varies source standing while holding contextual status fixed. Another holds the ordered events and local permissions fixed while varying whether \(R^{\mathrm{stip}}\) contains a cumulative or terminal condition that the trajectory violates. A matched placebo places the same restriction in quoted or otherwise non-operative text. Static transcript cases test local discrimination or reviewer reconstruction; a live tool-mediated harness is a separate, horizon-indexed projective target.

Purposive fixtures can expose level errors and local failure modes. They don't estimate real-world prevalence. These tests can show whether recorded classifications, observable decisions, predictions, or reviewer judgments respond to the tested contrasts. They don't by themselves establish internal representation, legal validity, or performance on other paths, systems, regimes, reviewer populations, deployments, or times.

\clearpage
\section{Conclusion}

Delegation assurance supplies regime-relative authorization semantics and typed representations for action and trajectory authorization, failure diagnosis, and comparative evaluation. These tools matter because AI deployments combine people, models, tools, persistent memory, and evaluators in ways that can obscure the distinctions among normative grants, recorded authorization states, proposed operations, executed trajectories, and assessment findings.

The framework treats the single action as the \(n=1\) trajectory case, represents local, prefix, cumulative, and terminal conditions, and distinguishes a normative hypergraph, recorded propagation graph, and execution trajectory. It also separates stipulated closed-world regimes from institutional regimes whose identity and applicability require adjudication. A record can support reconstruction without proving authority; an authorized action can be poorly recorded; and a pristine record can preserve an invalid grant. The executable artifact derives only the declared closed-world subset; it isn't an institutional-law engine.

What the framework establishes here is conceptual and instrumental rather than empirical. The semantics, the typed representations, and the separations they enforce stand on the arguments and the worked traces given above. The design analysis in Section~\ref{sec:comparative-test} is a further result of that kind: it shows by simulation that the gate forms an earlier version of this programme had frozen were defective in ways that no amount of data would have corrected, since conditioning a verdict on an interval lower bound selects inflated estimates whatever the sample. Retiring those gates didn't require the target studies, and neither does anything claimed in this paper.

Three separable tests would tell us something the arguments can't: randomized local discrimination, predictive representation ablation, and reviewer reconstruction ablation. Their specifications, input schemas, integrity rules, design floors, and analyzers were committed before any target outcome, which is itself a commitment rather than a promissory note, and only synthetic regression tests have been run against them. Until they are, the framework's usefulness for the questions those tests pose is undetermined, and no claim here rests on their results. Every extension beyond observed cases requires a projective declaration, intended use, warrant plan, and prospective revision rule. The framework has no scalar measure of authority, and noncompensation applies only where the regime or validated assessment use is conjunctive. Its scope remains limited to cases where action permissions, limits, authority transitions, trajectories, or authorization assessment are at issue; it doesn't turn every model error into a delegation failure.

\section*{Code and Data Availability}

\href{https://github.com/BrettRey/adversarial-pragmatics-for-ai-safety-evaluation}{The project repository} carries the delegation-assurance artifact under \path{assurance/delegation/}. The state described here is commit \texttt{95e927b}. A numbered release and an archival DOI are pending; until they exist, cite the commit rather than the branch.

The artifact contains three versioned schemas: \path{regime.schema.json} (\texttt{regime-2.0.0}), \path{authorization-trace.schema.json} (\texttt{authorization-trace-2.0.0}), and \path{oracle.schema.json} (\texttt{oracle-2.0.0}). It also contains the reviewer and predictor view policy, the frozen projective-claim register, seven fixtures encoding the level separations of Section~\ref{sec:typed-claim-stack}, a hidden-oracle set, and the frozen empirical program specifications with their input and output schemas under \path{assurance/delegation/empirical/}.

Two analyzers are included. The semantic engine derives authorization standing, guards, payload--scope fit, grant and subgrant validity, transition effects, local authorization, state continuity, prefix and cumulative conditions, terminal-state conditions, and the applicability tree from reference-free traces. The empirical analyzer validates the three frozen input schemas and emits observation vectors and noncompensatory family verdicts. Both are exercised by synthetic pass, fail, noncompensation, split, leakage, manifest, and oracle-masking regressions.

The hidden oracles are authored against the regime fixtures rather than generated from engine output, and the view policy excludes them from reviewer and predictor inputs. That construction detects disagreement between the engine and an independently written expectation; it doesn't establish that the two rest on different modelling assumptions, so it's weaker than a differential implementation and is reported as such.

No target-study outcomes exist. Every reported result is synthetic regression output. The engine derives the declared closed-world subset only, and doesn't implement open-text institutional interpretation, defeasible legal reasoning, arbitrary history predicates, or discretionary adjudication.


\clearpage
\begingroup
\renewcommand*{\bibfont}{\small}
\setlength{\bibitemsep}{0.25\baselineskip}
\printbibliography
\endgroup

\end{document}
